\section{定向与体积形式}\label{sec:orientation-and-volume-forms}

本节研究定向与微分形式之间的关系.

\begin{RiemannDefinition}\label{def:volume-form}
$n$ 维流形 $M$ 上的\textbf{体积形式} (或\textbf{体积元}) 是对所有 $p\in M$ 均有 $\omega_p\ne0$ 的 $n$-形式 $\omega$.
\end{RiemannDefinition}

体积形式的存在性等价于 $M$ 可定向.

\begin{RiemannProposition}\label{prop:orientation-and-existence-of-volume-form}
$n$ 维流形 $M$ 可定向, 当且仅当 $M$ 上存在体积形式.
\end{RiemannProposition}

\begin{RiemannProof}
设 $\omega$ 是 $M$ 上的体积形式, 考虑图册 $\{(U_\alpha,\varphi_\alpha)\}$. 不失一般性, 可假设开集 $U_\alpha$ 连通. 我们将据此构造一个新图册, 使其重叠映射的导数具有正的行列式. 事实上, 考虑 $\omega$ 在其中一个开集 $U_\alpha\subset\mathbb{R}^n$ 上的表示, 有
\[
  \varphi_\alpha^*\omega=a_\alpha\,\mathrm{d}x_\alpha^1\wedge\cdots\wedge\mathrm{d}x_\alpha^n,
\]
其中函数 $a_\alpha$ 不能为零, 因而符号固定. 若 $a_\alpha$ 为正, 保留相应的参数化. 否则, 将 $\varphi_\alpha$ 与下列映射复合, 构造一个新的参数化:
\[
  (x^1,\ldots,x^n)\mapsto(-x^1,x^2,\ldots,x^n).
\]
在这些新坐标下, 新的函数 $a_\alpha$ 显然为正. 对所有坐标邻域重复这个过程, 得到一个所有函数 $a_\alpha$ 均为正的新图册, 仍记为 $\{(U_\alpha,\varphi_\alpha)\}$. 此外, 只要 $W:=\varphi_\alpha(U_\alpha)\cap\varphi_\beta(U_\beta)\ne\varnothing$, 就有 $\omega_\alpha=(\varphi_\beta^{-1}\circ\varphi_\alpha)^*\omega_\beta$. 因此,
\begin{align*}
  a_\alpha\,\mathrm{d}x_\alpha^1\wedge\cdots\wedge\mathrm{d}x_\alpha^n
  &=(\varphi_\beta^{-1}\circ\varphi_\alpha)^*
    (a_\beta\,\mathrm{d}x_\beta^1\wedge\cdots\wedge\mathrm{d}x_\beta^n)\\
  &=(a_\beta\circ\varphi_\beta^{-1}\circ\varphi_\alpha)
    \bigl(\det(\mathrm{d}(\varphi_\beta^{-1}\circ\varphi_\alpha))\bigr)
    \,\mathrm{d}x_\alpha^1\wedge\cdots\wedge\mathrm{d}x_\alpha^n
\end{align*}
从而 $\det(\mathrm{d}(\varphi_\beta^{-1}\circ\varphi_\alpha))>0$. 因此 $M$ 可定向.

反过来, 若 $M$ 可定向, 则考虑一个图册 $\{(U_\alpha,\varphi_\alpha)\}$, 其重叠映射 $\varphi_\beta^{-1}\circ\varphi_\alpha$ 满足 $\det\mathrm{d}(\varphi_\beta^{-1}\circ\varphi_\alpha)>0$. 取从属于 $M$ 的相应坐标邻域覆盖的单位分解 $\{\rho_i\}_{i\in I}$, 可定义形式
\[
  \omega_i:=\rho_i\,\mathrm{d}x_i^1\wedge\cdots\wedge\mathrm{d}x_i^n
\]
使 $\operatorname{supp}\omega_i=\operatorname{supp}\rho_i\subset\varphi_{\alpha_i}(U_{\alpha_i})$. 令这些形式在 $\operatorname{supp}\rho_i$ 外为零, 从而将其延拓到 $M$ 上, 可定义形式 $\omega:=\sum_{i\in I}\omega_i$. 显然, $\omega$ 是 $M$ 上一个良定义的 $n$-形式, 所以只需证明对所有 $p\in M$, 都有 $\omega_p\ne0$. 设 $p$ 是 $M$ 中一点. 存在 $i\in I$ 使 $\rho_i(p)>0$, 因而存在线性无关的向量 $v_1,\ldots,v_n\in T_pM$, 使 $(\omega_i)_p(v_1,\ldots,v_n)>0$. 此外, 对所有其余的 $j\in I\setminus\{i\}$, 都有 $(\omega_j)_p(v_1,\ldots,v_n)\geq0$. 事实上, 若 $p\notin\varphi_{\alpha_j}(U_{\alpha_j})$, 则 $(\omega_j)_p(v_1,\ldots,v_n)=0$. 另一方面, 若 $p\in\varphi_{\alpha_j}(U_{\alpha_j})$, 则由 \eqref{eq:pullback-top-form-in-coordinates},
\[
  \mathrm{d}x_j^1\wedge\cdots\wedge\mathrm{d}x_j^n
  =\det(\mathrm{d}(\varphi_{\alpha_j}^{-1}\circ\varphi_{\alpha_i}))
    \,\mathrm{d}x_i^1\wedge\cdots\wedge\mathrm{d}x_i^n
\]
因而
\[
  (\omega_j)_p(v_1,\ldots,v_n)
  =\frac{\rho_j(p)}{\rho_i(p)}
    \bigl(\det(\mathrm{d}(\varphi_{\alpha_j}^{-1}\circ\varphi_{\alpha_i}))\bigr)
    (\omega_i)_p(v_1,\ldots,v_n)\geq0.
\]
因此 $\omega_p(v_1,\ldots,v_n)>0$, 所以 $\omega$ 是体积形式.
\end{RiemannProof}

\begin{RiemannRemark}\label{rem:orientation-determined-by-volume-form}
有时把体积形式称为一个定向. 在这种情况下, $M$ 上的定向是这样的: $T_pM$ 的一组基 $\{v_1,\ldots,v_n\}$ 为正定向的, 当且仅当 $\omega_p(v_1,\ldots,v_n)>0$.
\end{RiemannRemark}

在可定向流形 $M$ 上固定一个体积形式 $\omega\in\Omega^n(M)$ 后, 可以定义任意具有紧支集的函数 $f\in C^\infty(M,\mathbb{R})$ 的\textbf{积分}为
\[
  \int_Mf:=\int_Mf\omega
\]
(其中 $M$ 的定向由 $\omega$ 决定). 若 $M$ 紧, 定义其\textbf{体积}为
\[
  \operatorname{vol}(M):=\int_M1=\int_M\omega.
\]

\begin{RiemannExercise}\label{ex:orientation-and-volume-forms}
\begin{enumerate}[label=(\arabic*),ref=\arabic*,leftmargin=*]
\item\label{ex:orientability-of-product} 证明 $M\times N$ 可定向当且仅当 $M$ 和 $N$ 都可定向.
\item\label{ex:positive-integral-against-volume} 设 $M$ 是紧定向流形, 体积元为 $\omega\in\Omega^n(M)$. 证明若 $f>0$, 则 $\int_Mf\omega>0$.
{\small (\textbf{Remark:} 特别地, 紧流形的体积总是正的).}
\item\label{ex:vanishing-exterior-derivative-on-compact-orientable} 设 $M$ 是 $n$ 维紧可定向流形, $\omega$ 是 $M$ 中的 $(n-1)$-形式.
\begin{enumerate}[label=(\alph*),leftmargin=*]
\item 证明存在一点 $p\in M$, 使得 $(\mathrm{d}\omega)_p=0$.
\item 证明不存在从单位圆到 $\mathbb{R}$ 的浸入 $f:S^1\to\mathbb{R}$.
\end{enumerate}
\item\label{ex:forms-and-orientation-on-projective-space} 设 $f:S^n\to S^n$ 是对径映射. 回顾 $n$ 维射影空间是可微流形 $\mathbb{R}P^n=S^n/\mathbb{Z}_2$, 其中群 $\mathbb{Z}_2=\{1,-1\}$ 通过 $1\cdot x=x$ 和 $(-1)\cdot x=f(x)$ 作用于 $S^n$. 设 $\pi:S^n\to\mathbb{R}P^n$ 是自然投影.
\begin{enumerate}[label=(\alph*),leftmargin=*]
\item 证明 $\omega\in\Omega^k(S^n)$ 可写成 $\omega=\pi^*\theta$, 其中 $\theta\in\Omega^k(\mathbb{R}P^n)$, 当且仅当 $f^*\omega=\omega$.
\item 证明 $\mathbb{R}P^n$ 可定向当且仅当 $n$ 为奇数, 且在此情形下,
\[
  \int_{S^n}\pi^*\theta=2\int_{\mathbb{R}P^n}\theta.
\]
\item 证明当 $n$ 为偶数时, 球面 $S^n$ 是 $\mathbb{R}P^n$ 的可定向二重覆盖 [参见第~\ref{chap:differentiable-manifolds}~章的 Exercise~\ref{ex:orientability}(\ref{ex:orientable-double-covering})].
\end{enumerate}
\item\label{ex:divergence-theorem-for-volume-form} 设 $M$ 是紧定向带边界流形, $\omega\in\Omega^n(M)$ 是体积元. 向量场 $X\in\mathfrak{X}(M)$ 的\textbf{散度}是满足
\[
  L_X\omega=(\operatorname{div}(X))\omega
\]
的函数 $\operatorname{div}(X)$ [参见 Exercise~\ref{ex:differential-forms}(\ref{ex:lie-derivative-of-differential-form})]. 证明
\[
  \int_M\operatorname{div}(X)=\int_{\partial M}\iota(X)\omega.
\]
\item\label{ex:brouwer-fixed-point-theorem} (\emph{Brouwer 不动点定理})
\begin{enumerate}[label=(\alph*),leftmargin=*]
\item 设 $M$ 是 $n$ 维紧可定向带边界流形, 且 $\partial M\ne\varnothing$. 证明不存在满足 $f|_{\partial M}=\operatorname{id}$ 的光滑映射 $f:M\to\partial M$.
\item 证明 \textbf{Brouwer 不动点定理}: 闭球 $B:=\{x\in\mathbb{R}^n\mid\lVert x\rVert\leq1\}$ 到自身的任意光滑映射 $g:B\to B$ 都有一个不动点, 即存在一点 $p\in B$ 使得 $g(p)=p$.
{\small (\textbf{提示:} 对每个点 $x\in B$, 考虑从 $g(x)$ 出发且经过 $x$ 的射线 $r_x$. 在 $r_x\cap\partial B$ 上恰有一点 $f(x)$ 不同于 $g(x)$. 考虑映射 $f:B\to\partial B$).}
\end{enumerate}
\end{enumerate}
\end{RiemannExercise}
